\section{Finite random routing and exact conditional failure}
\label{sec:routing}

This section presents the routing argument inherited from
\cite[Sections~3--4]{OpenAI84}, with the common grids and sampled terms
from \cref{sec:synchronization}. We include the conditioning details
because the points choose their leaves adaptively: unconditional
independence of their entire routes is neither true nor needed.

\subsection{Selectors, defaults, and terminal bits}

Set $m_0=M-1$. For every edge $e=(P,i)$ with $1\le i<M$, assign an
independent fair bit to every key at resolution $b_e$:
\[
 S_e:\{0,\ldots,N_{b_e}-1\}\longrightarrow\{0,1\}.
\]
The $M$th child is the default and has no selector table. For each
leaf $L$, let $b(L)$ be the endpoint of its incoming edge, and assign
independent Bernoulli-$p$ entries to a terminal table $T_L$ on the
keys at resolution $b(L)$. Every entry in all these tables is
independent of all others. The probability space is finite.

Route $z\in\R$ from the root. At each internal vertex choose the
first child $i<M$ for which $S_{(P,i)}(J_{b_{(P,i)}}(z))=1$;
if no such child exists, choose child $M$. Write $L(z)$ for the leaf
reached. Define the random set
\begin{equation}\label{eq:random-set}
 A=\{z:T_{L(z)}(J_{b(L(z))}(z))=1\}.
\end{equation}
Each realization is a $1$-periodic finite union of half-open cells.
All its tables and routing rules are independent of $q$ and $t$.

\begin{lemma}[Mean density]\label{lem:mean-density}
$\E\rho(A)=p$.
\end{lemma}

\begin{proof}
Let $\calA$ be the sigma-field of all selector entries. Conditional
on $\calA$, the terminal address of each fixed point $z$ is determined.
Its terminal bit remains Bernoulli-$p$, so
$\Prob(z\in A\mid\calA)=p$. Integrate on one period; since the
probability space is finite, the interchange of expectation and
integration is immediate.
\end{proof}

\subsection{Expose only the center entries}

For a fixed center $x$, reveal the addressed entry at $x$ in \emph{every}
selector table, including tables outside its route. Define
\begin{equation}\label{eq:center-field}
 \calF_x=\sigma\bigl(S_e(J_{b_e}(x)): e=(P,i),\ i<M\bigr)
 \subseteq\calA.
\end{equation}
The exposed addresses are deterministic. On any atom $\xi$ of $\calF_x$,
all other selector entries remain independent fair bits and all terminal
entries retain their independent Bernoulli-$p$ distributions.

At each vertex reached by $x$, its $m_0$ selector entries are fresh
relative to the previous ancestor decisions. Thus
\begin{equation}\label{eq:no-default}
 \Prob(\text{$x$ has no default on its route})=(1-2^{-m_0})^d.
\end{equation}
On an atom where the route uses a default, let $U$ be the first such
vertex and $h$ its height. Both are fixed on that atom. Put $r=r_h$
and $e_i=(U,i)$ for $1\le i\le m_0$. All selectors at $x$ on these
edges are zero.

For $z\in\R$, let $L_i(z)$ be the leaf reached by routing from
child $U_i$. Define a local test
\begin{equation}\label{eq:local-test}
 Q_i(z)=S_{e_i}(J_{b_{e_i}}(z))\,
 T_{L_i(z)}(J_{b(L_i(z))}(z)).
\end{equation}
This uses only the table on $e_i$ and tables below that child.

\begin{lemma}[A successful local test is an actual hit]\label{lem:local-hit}
Suppose $x\in\calS$ and the atom $\xi$ has first default $U$.
For $j\in W_{e_i}$ and any $q\in[\alpha,\beta]$, $t\in[1,2]$,
\begin{equation}\label{eq:local-hit}
 Q_i(x+ta_j(q))=1\quad\Longrightarrow\quad x+ta_j(q)\in A.
\end{equation}
\end{lemma}

\begin{proof}
At a strict ancestor of $U$, the child chosen by $x$ is nondefault.
Every selector consulted to make that choice lies earlier than $e_i$
in preorder. By \cref{lem:stable}, its value at $x+ta_j(q)$ equals
its value at $x$. Hence the translated point reaches $U$.
The selectors of children $1,\ldots,i-1$ of $U$ are also earlier
in preorder and remain zero. If the test in \eqref{eq:local-test}
is one, child $i$ is selected and its terminal entry is one. This is
exactly the membership rule \eqref{eq:random-set}.
\end{proof}

\begin{lemma}[Exact conditional failure]\label{lem:failure}
Under the assumptions of \cref{lem:local-hit}, for every fixed
$(q,t)\in[\alpha,\beta]\times[1,2]$,
\begin{equation}\label{eq:failure}
 \Prob\!\left(Q_i(x+ta_j(q))=0\text{ for every }i\le m_0,
 j\in W_{e_i}\,\middle|\,\xi\right)
 =(1-p/2)^{m_0r}.
\end{equation}
\end{lemma}

\begin{proof}
There are $m_0r$ pairs $(i,j)$. Write
\[
 A_{ij}=S_{e_i}(J_{b_{e_i}}(x+ta_j(q))).
\]
For one $i$, \cref{lem:separation} makes all these addresses distinct
and different from the exposed center address. Different $i$ use
different tables. Conditional on $\xi$, the $A_{ij}$ are therefore
independent fair bits.

Now condition on the larger sigma-field $\calA$ of \emph{all}
selectors. The leaf and terminal address of every local test are fixed.
Those addresses are pairwise distinct. Different children $i$ have
disjoint subtrees and hence different terminal tables. Within a child,
different leaves again have different tables. If two distinct $j,j'$
reach the same leaf, its resolution is at least $b_{e_i}$, so
\cref{lem:separation} makes their terminal keys distinct. This covers
all possible pairs.

Consequently, conditional on $\calA$, the probability that all tests
fail is
\[
 (1-p)^{\sum_{i,j}A_{ij}}.
\]
The tests with $A_{ij}=0$ already fail; each of the others consults a
different Bernoulli-$p$ terminal entry. Average this expression over
the independent fair $A_{ij}$ conditional on $\xi$:
\[
 \E\!\left[(1-p)^{\sum_{i,j}A_{ij}}\mid\xi\right]
 =\prod_{i,j}\left(\tfrac12+\tfrac12(1-p)\right)
 =(1-p/2)^{m_0r}.
\]
No independence of the internal routes has been invoked.
\end{proof}

The geometry now enters only through how many different vectors of
local keys can occur as $(q,t)$ varies. That count must be obtained
\emph{before} revealing the unexposed table entries; otherwise a
union bound over representatives would not be justified.
